% Драфт заполнения characteristic.tex (Общая характеристика работы).
% Источник содержания: CORE.md (центральная идея, ключевые результаты, вклад автора).
% Статус: черновик для ревью. В common/characteristic.tex НЕ внесён.
%
% ОБЛАСТЬ ПРАВКИ: заполнены пустые \TODO-блоки — \aim, \tasks, \novelty,
% \influence, \methods, \defpositions, \reliability, \contribution.
% НЕ МЕНЯЮТСЯ: актуальность (\actuality, вверху файла), апробация (\probation,
% список конференций) и блок публикаций (refsection) — остаются как есть.
%
% Места, требующие подтверждения автора, помечены: % [ПОДТВЕРДИТЬ].

% ============================================================
% [сохраняется без изменений: \actuality — актуальность работы]
% ============================================================


Несмотря на обилие методов анализа нейронной активности, они, как правило,
работают на одном масштабе: либо описывают популяцию как единое целое (методы
снижения размерности, модели латентной динамики), либо характеризуют
селективность отдельных нейронов, --- и редко связывают эти уровни между собой.
Между тем коллективная динамика нейронной активности обладает низкоразмерной
структурой, которую продуктивно изучать одновременно на двух уровнях:
популяционном (скрытые переменные, геометрия многообразий, размерность) и
индивидуальном (селективность отдельных нейронов к поведенческим и латентным
переменным). Более того, эти подходы не обязаны быть привязаны к биологическому
субстрату: один и тот же аналитический инструментарий применим и к нейронам
мозга, и к элементам искусственных нейронных сетей, что открывает путь к
переносу методов между нейронаукой и машинным обучением. Настоящая работа
развивает методы анализа скрытых переменных коллективной динамики на обоих
уровнях и демонстрирует их субстрат-независимость.

{\aim} данной работы является разработка методов выявления и количественной
характеризации скрытых переменных коллективной динамики нейронной активности
на популяционном и индивидуальном уровнях, а также их применение к
биологическим и искусственным нейронным сетям.

Для~достижения поставленной цели необходимо было решить следующие {\tasks}:
\begin{enumerate}[beginpenalty=10000]
  \item Адаптировать нелинейные методы снижения размерности для восстановления
  геометрии пространства состояний из популяционной активности нейронов
  гиппокампа и определить границы их применимости в сравнении с линейными
  методами.
  \item Разработать алгоритм оценки эффективной размерности нейронной
  активности, учитывающий конечность экспериментальных записей, и исследовать
  связь эффективной и внутренней размерности с поведением животного.
  \item Разработать метод анализа коллективной динамики, явно учитывающий
  темпоральную структуру активности, --- мультиплексный граф рекуррентности
  популяции --- и апробировать его на данных кальциевого имиджинга.
  \item Создать информационно-теоретический метод выявления селективности
  отдельных нейронов, обеспечивающий разделение вкладов коррелированных
  переменных, и применить его к биологическим и искусственным нейронным сетям.
  \item Разработать безградиентные методы синтеза оптимальных стимулов
  (максимизации активации) для нейронов искусственных, в том числе импульсных,
  нейронных сетей.
  \item Объединить разработанные методы анализа в единую программную среду
  с инструментами автоматизированной предобработки экспериментальных данных
  кальциевого имиджинга.
\end{enumerate}


{\novelty}
\begin{enumerate}[beginpenalty=10000]
  \item Впервые алгоритм последовательной коррекции спектра корреляционных
  матриц применён к экспериментальным записям нейронной активности; показано,
  что эффективная размерность снижается в периоды движения животного.
  \item Впервые для нелинейной меры размерности (внутренней) обнаружено
  степенное масштабирование с числом регистрируемых нейронов.
  \item Предложен мультиплексный граф рекуррентности популяции --- оригинальный
  способ агрегации индивидуальных графов рекуррентности в единое представление;
  впервые метод применён к сырым сигналам кальциевого имиджинга \textit{in vivo}.
  \item Показано, что значительная доля наблюдаемой мультиселективности
  нейронов является артефактом коррелированности поведенческих переменных,
  а замена взаимной информации линейной корреляцией упускает большую часть
  значимых связей.
  \item Впервые реализован безградиентный синтез оптимальных стимулов
  (максимизация активации) для импульсных нейронных сетей.
  \item Разработан программный фреймворк, впервые объединяющий индивидуальный,
  популяционный и сетевой масштабы анализа нейронной активности в единой среде.
  \item Показано, что сочетание автокодировщика и метода анализа селективности
  вскрывает вычислительный механизм рекуррентной нейронной сети: организацию
  пространства состояний (бассейны притяжения) и качественно различные
  вычислительные стратегии, возникающие при разных алгоритмах обучения при
  одинаковой точности.
  \item Продемонстрирована субстрат-независимость разработанных методов:
  их применимость к биологическим и искусственным (рекуррентным и импульсным)
  нейронным сетям.
\end{enumerate}


{\influence} Разработанные методы и открытый программный пакет применимы
к широкому классу данных --- записям кальциевого имиджинга, электрофизиологии
и активациям искусственных нейронных сетей --- и предоставляют единый
инструментарий анализа на индивидуальном, популяционном и сетевом масштабах.
Количественные характеристики размерности и коллективной динамики дают новые
инструменты для изучения популяционного кода, а методы анализа индивидуальной
селективности и синтеза оптимальных стимулов --- для интерпретации функции
отдельных нейронов. Инструменты автоматизированной предобработки существенно
сокращают трудозатраты при анализе записей кальциевого имиджинга.


{\methods} В работе использованы: кальциевый имиджинг \textit{in vivo}
с последующей предобработкой (декомпозиция методом CNMF-E, автоматизированный
контроль качества выделенных компонент); нелинейные методы снижения
размерности (лапласовы собственные карты, Isomap, UMAP, автокодировщики);
методы оценки размерности (коррекция спектра корреляционных матриц, оценщик
2-NN); методы нелинейной динамики (вложение Такенса, графы рекуррентности,
рекуррентный количественный анализ); теоретико-информационные методы
(взаимная информация, оценщик на основе копулы GCMI, перестановочное
тестирование с поправками на множественные сравнения); безградиентная
оптимизация на основе разложения тензорного поезда; статистический анализ
(дисперсионный анализ, суррогатные данные).


{\defpositions}
% [ПОДТВЕРДИТЬ] число и форму положений желательно сверить с правилами
% диссертационного совета.
\begin{enumerate}[beginpenalty=10000]
  \item Нелинейные методы снижения размерности и мультиплексный граф
  рекуррентности популяции восстанавливают геометрию пространства состояний
  (геометрию исследуемой среды) из коллективной активности нейронов гиппокампа,
  недоступную линейным методам; единое представление мультиплексного графа
  одновременно восстанавливает геометрию среды и разделяет динамические режимы
  коллективной активности.
  \item Эффективная и внутренняя размерности количественно характеризуют
  многообразие нейронной активности: внутренняя размерность существенно ниже
  эффективной и демонстрирует степенное масштабирование с числом нейронов,
  а эффективная размерность обратно связана со скоростью движения животного.
  \item Разработанный безградиентный метод синтеза оптимальных стимулов
  (максимизации активации) служит инструментом изучения внутренних представлений
  искусственных, в том числе импульсных, нейронных сетей и связывает
  индивидуальную селективность отдельных нейронов с популяционными
  характеристиками их активности (динамикой размерности и информации по слоям
  сети).
  \item Создан единый программный фреймворк, объединяющий анализ индивидуальной
  селективности нейронов (INTENSE), популяционной динамики и сетевой структуры;
  разработанные методы субстрат-независимы и применимы как к биологическим, так
  и к искусственным, в том числе импульсным, нейронным сетям.
\end{enumerate}


{\reliability} полученных результатов обеспечивается применением статистических
методов с контролем значимости (перестановочное тестирование с поправками
на множественные сравнения, суррогатные данные); валидацией разработанных
методов на синтетических данных с аналитически известными свойствами;
согласованностью результатов между экспериментальными сессиями и животными;
воспроизводимостью вычислений за счёт открытого программного кода; апробацией
результатов на международных научных конференциях и публикациями
в рецензируемых научных журналах.


% ============================================================
% [сохраняется без изменений: \probation — список конференций]
% ============================================================


{\contribution} Автором лично разработаны методы оценки эффективной
и внутренней размерности нейронной активности, информационно-теоретический
метод анализа селективности нейронов INTENSE, метод мультиплексного сетевого
анализа коллективной динамики, программный пакет DRIADA и инструменты
автоматизированного контроля качества данных кальциевого имиджинга. Автору
принадлежат: анализ данных в исследовании по восстановлению геометрии среды
(экспериментальные записи получены коллегами); применение автокодировщиков
и метода INTENSE к рекуррентным нейронным сетям. Весь анализ селективности
и внутренних представлений нейронов импульсных сетей выполнен автором;
программная система MANGO для безградиентного синтеза оптимальных стимулов
создавалась совместно с соавторами. Экспериментальные записи нейронной активности
получены коллегами.


% ============================================================
% [сохраняется без изменений: блок \publications (refsection)]
% ============================================================
